% -----------------------------------------------------------------------------
% This file is automatically generated.
% Do not edit manually.
% Source: notebooks/support/regression-analysis/support.Rmd
% -----------------------------------------------------------------------------


\noindent The \ttblue{summary} command produces summary statistics of all variables in the dataset, the length (number of data points) for character types (such as \ttblue{month}) and distribution statistics for numeric types.

\par\addvspace{\topsep}
\begingroup
\fvset{listparameters={%
  \setlength{\topsep}{0pt}%
  \setlength{\partopsep}{0pt}%
  \setlength{\parsep}{0pt}%
  \setlength{\itemsep}{0pt}%
}}
\begin{Verbatim}[breaklines=true,formatcom=\color{ada_blue}]
summary(USSteamCo)
\end{Verbatim}
\begin{Verbatim}[breaklines=true,formatcom=\color{ada_light_blue}]
    month              revenue           production
 Length:48          Min.   : 4809952   Min.   :151104
 Class :character   1st Qu.:10325796   1st Qu.:371474
 Mode  :character   Median :13542764   Median :502280
                    Mean   :14032843   Mean   :493034
                    3rd Qu.:17865242   3rd Qu.:646926
                    Max.   :26792280   Max.   :805600
     coolDD          heatDD
 Min.   :  0.0   Min.   :   0.0
 1st Qu.:  0.0   1st Qu.:  10.5
 Median : 21.0   Median : 258.5
 Mean   :113.3   Mean   : 367.3
 3rd Qu.:203.8   3rd Qu.: 690.0
 Max.   :505.0   Max.   :1123.0
\end{Verbatim}
\endgroup
\par\addvspace{\topsep}

\noindent We create a dummy variable to distinguish between heating months and cooling months. Such a variable would have the value 1 in heating months and 0 in cooling months. The effect of this variable on the trajectory of the regression line is explained in Section 5.4.4.
